\section{Lightweight Iterative Optimization Method}
\label{sec:liom}
With the flight horizon fixed, LIOM revises the trajectories and their local collision-avoidance regions together. Its corridor reconstruction and external-penalty principles build on~\cite{bu2025,li2022liom}. The procedure developed here organizes computation into a hard-QP entry and two recovery layers. The first layer relaxes conflicting kinematic and relative-corridor restrictions so that geometry can be updated. The second layer responds to a reconstructed motion that violates a physical limit, using tighter recovery and an exact-kinematics finish. Three outcomes are distinguished throughout: a nonempty relaxed subproblem, numerical recovery of its candidate, and final acceptance of the reconstructed continuous motion.

\subsection{Two-Penalty Recovery Formulation}
For the double-integrator motion model, the kinematic defect is affine in the position and velocity variables. The two quantities that measure constraint recovery are
\begin{equation}
 D=\sum_{i=1}^N\sum_{k=0}^{K-1}\norm{\mathbf E_{i,k}}_2^2,\qquad
 S=\sum_{(i,j,k)\in\mathcal A}\delta_{ij,k}.
 \label{eq:recovery}
\end{equation}
At iteration $\nu$, the augmented objective is
\begin{equation}
 J^{(\nu)}=C+\frac{\rho_q^{(\nu)}}hD+\rho_s^{(\nu)}S.
 \label{eq:penalty}
\end{equation}
The quantities $D$ and $S$ have units of $\mathrm{m}^2$ and $\mathrm{m}$, respectively. The factor $1/h$ accounts for interval length in the kinematic penalty. For a dimensionless augmented objective in SI coordinates, $\rho_q$ and $\rho_s$ have units $\mathrm{s}/\mathrm{m}^2$ and $\mathrm{m}^{-1}$, respectively. Multiplying the entire objective by a positive scalar improves numerical scaling without changing its minimizers. The weights $\rho_q^{(\nu)}$ and $\rho_s^{(\nu)}$ regulate the recovery of kinematics and relative-corridor compatibility, respectively.

Let $\mathcal I=\{1,\ldots,N\}$, $\mathcal K=\{0,\ldots,K\}$ and $\mathcal K^-=\{0,\ldots,K-1\}$. Corridor bounds and $\mathcal A^{(\nu)}$ are fixed during an inner solve. The complete subproblem is
\begin{adjustwidth}{-\extralength}{0cm}
\begin{subequations}\label{eq:subproblem}\begin{align}
 \min_{\mathbf P,\mathbf V,\boldsymbol\delta}\quad&
 C(\mathbf V)
 +\frac{\rho_q^{(\nu)}}h\sum_{i\in\mathcal I}\sum_{k\in\mathcal K^-}\norm{\mathbf E_{i,k}}_2^2
 +\rho_s^{(\nu)}\sum_{(i,j,k)\in\mathcal A^{(\nu)}}\delta_{ij,k}
 \label{eq:nlpobjective}\\
 \mathrm{s.t.}\quad&
 \underline{\mathbf w}+r\one\leq\bp_{i,k}\leq\overline{\mathbf w}-r\one,
 && i\in\mathcal I,\ k\in\mathcal K,\label{eq:nlpworkspace}\\
 &\underline{\boldsymbol\gamma}_{i,k}\leq\bp_{i,k}\leq\overline{\boldsymbol\gamma}_{i,k},
 &&i\in\mathcal I,\ k\in\mathcal K,\label{eq:nlpstc}\\
 &U\bv_{i,k}\leq v_{\max}\mathbf b,
 &&i\in\mathcal I,\ k\in\mathcal K,\label{eq:nlpvelocity}\\
 &U(\bv_{i,k+1}-\bv_{i,k})/h\leq a_{\max}\mathbf b,
 &&i\in\mathcal I,\ k\in\mathcal K^-,\label{eq:nlpacceleration}\\
 &\bp_{i,0}=\mathbf s_i,\quad\bp_{i,K}=\mathbf g_i,\quad
 \bv_{i,0}=\bv_{i,K}=\mathbf0,
 &&i\in\mathcal I,\label{eq:nlpboundary}\\
 &\underline{\boldsymbol\eta}_{ij,k}-\delta_{ij,k}\one
 \leq\bp_{i,k}-\bp_{j,k},
 &&(i,j,k)\in\mathcal A^{(\nu)},\label{eq:nlprellower}\\
 &\bp_{i,k}-\bp_{j,k}\leq\overline{\boldsymbol\eta}_{ij,k}+\delta_{ij,k}\one,
 &&(i,j,k)\in\mathcal A^{(\nu)},\label{eq:nlprelupper}\\
 &\delta_{ij,k}\geq0,
 &&(i,j,k)\in\mathcal A^{(\nu)}.\label{eq:nlpslack}
\end{align}\end{subequations}
\end{adjustwidth}
In the objective, $i$ ranges over $\mathcal I$. Acceleration and $\mathbf E_{i,k}$ are substituted expressions. Workspace and STC bounds are intersected, and endpoint states are imposed by equal lower and upper variable bounds. Every constraint is linear or a variable bound. Since $C,D$ are convex quadratic and $S$ is linear, the fixed-corridor program is a convex quadratic program (QP). The affine double-integrator equations and hard relative boxes also retain this convex structure. The two relaxations preserve availability of intermediate feasible points when those hard requirements conflict; corridor reconstruction changes the local geometric region between solves.

\begin{Proposition}[Nonemptiness of the relaxed fixed-corridor subproblem]
\label{prop:nonempty}
Suppose every spatial box is nonempty, contained in the admissible workspace, and contains the prescribed start or goal at the corresponding endpoint. For the rest-to-rest task, the relaxed subproblem~\eqref{eq:subproblem} has a feasible point with finite objective.
\end{Proposition}
\begin{proof}
Choose any position inside each box, preserve the fixed endpoint positions, and set every velocity to zero. The velocity, acceleration, and rest-to-rest boundary bounds are satisfied. For each active triple, choose
\begin{equation}
 \delta_{ij,k}=\max\!\left\{0,\,
 \max_d(\underline{\eta}_{ij,k,d}-r_{ij,k,d}),\,
 \max_d(r_{ij,k,d}-\overline{\eta}_{ij,k,d})\right\}.
 \label{eq:feasiblepoint}
\end{equation}
Both relative inequalities are satisfied. All chosen states and slacks are finite, hence so is the augmented objective.
\end{proof}
Proposition~\ref{prop:nonempty} establishes local subproblem nonemptiness under the stated box conditions. The outer loop checks numerical recovery and continuous-motion acceptance separately. A failed construction or an exhausted budget terminates the attempt; a finite candidate can support the next corridor update even when the current solve has not reached numerical recovery.

\subsection{Hard-QP Entry and Adaptive Recovery}
An initial hard QP provides a fast path when the reference already admits a consistent motion. It retains the spatial boxes and motion bounds, imposes $\mathbf E_{i,k}=\mathbf0$, and replaces the relative boxes by separating half-spaces
\begin{equation}
 \mathbf n_{ij,k}^{\top}(\bp_{i,k}-\bp_{j,k})\geq R_{ij,k},\qquad
 \mathbf n_{ij,k}=\frac{\widehat{\bp}_{i,k}-\widehat{\bp}_{j,k}}
 {\norm{\widehat{\bp}_{i,k}-\widehat{\bp}_{j,k}}_2}.
 \label{eq:hardentry}
\end{equation}
Coincident reference nodes use the assigned start separation to define a direction. This entry solve receives a bounded fraction of the planning budget. A candidate that passes reconstruction and complete-interval checks is returned immediately. Otherwise, the two-penalty box formulation supplies the recovery loop. The fallback is triggered by the outcome of the hard solve and its motion check.

For the relaxed loop, normalize the recovery quantities as
\begin{equation}
 \bar D=\frac{D}{NK}\left(\frac{h_0}{h}\right)^2,\qquad
 \bar S=\frac{S}{\max\{1,|\mathcal A|\}},\qquad
 \bar D\leq\varepsilon_q,\quad\bar S\leq\varepsilon_s.
 \label{eq:tolerances}
\end{equation}
Here $h_0$ is a reference interval length. The defect test measures the mean squared displacement mismatch at a common temporal scale; the slack test measures mean relative-box violation. Their units are $\mathrm{m}^2$ and $\mathrm{m}$. This normalization removes direct summation over team size and active pairs. Final physical acceptance remains a separate, unnormalized test of every returned interval.
Each penalty is increased only when its corresponding recovery test fails:
\begin{equation}
 \rho_a^{(\nu+1)}=
 \begin{cases}
 \min\{\beta_a\rho_a^{(\nu)},\overline\rho_a\},&Q_a^{(\nu)}>\varepsilon_a,\\
 \rho_a^{(\nu)},&Q_a^{(\nu)}\leq\varepsilon_a,
 \end{cases}
 \quad (Q_q,Q_s)=(\bar D,\bar S).
 \label{eq:weights}
\end{equation}
Capping limits the penalty scale while reconstruction changes the local feasible region. A finite bounded candidate seeds the next corridors. If the solver provides no usable iterate, the spatial-box centers provide a shared relocation rule, with fixed endpoints preserved. Static guides remain available to relocate obstacle-intersecting seeds. Hard iterative comparisons receive the same updates and planning budget.

Numerical recovery and a bounded violation of the local linear constraints trigger control reconstruction. If the reconstructed motion fails physical acceptance, the method tightens $\varepsilon_q$, increases the kinematic penalty within its cap, and attempts an exact-kinematics, zero-slack QP in the current boxes. A successful motion is returned; otherwise the relaxed candidate seeds another corridor update. Thus an unsuccessful physical check has an explicit continuation branch. Exhausting the iteration or computation budget returns planning failure.

For a feasible fixed convex problem, a linear slack penalty can be exact above a threshold determined by suitable optimal constraint multipliers. That threshold varies with the current corridors. A quadratic defect penalty generally retains nonzero defects at finite weight. Capped continuation, exact finishing, and complete-interval checks jointly determine the accepted motion. Acceleration effort regularizes each local solution.

\subsection{Exact Rest-to-Rest Output from Recovered Controls}
\label{sec:arcs}
Finite kinematic defects are handled directly at the control level. For one spatial coordinate of UAV $i$, collect candidate accelerations in $\mathbf a_i\in\R^K$. In the following two equations, $s_i,g_i,v_i^0,v_i^f$ denote components along this same coordinate. Constant-acceleration integration expresses terminal velocity and displacement as two control moments:
\begin{equation}
 H=
 \begin{bmatrix}
 h&h&\cdots&h\\
 h^2(K-\tfrac12)&h^2(K-\tfrac32)&\cdots&h^2/2
 \end{bmatrix},\qquad
 \mathbf b_i=\begin{bmatrix}v_i^f-v_i^0\\g_i-s_i-T_f v_i^0\end{bmatrix}.
 \label{eq:moments}
\end{equation}
The minimum Euclidean-norm correction satisfying $H\bar{\mathbf a}_i=\mathbf b_i$ is
\begin{equation}
 \bar{\mathbf a}_i=\mathbf a_i+
 H^\top(HH^\top)^{-1}(\mathbf b_i-H\mathbf a_i).
 \label{eq:correction}
\end{equation}
The two rows of $H$ are independent for $K\geq2$, so only a $2\times2$ linear system is required for each coordinate. The correction requires only a closed-form linear solve.

Starting from the prescribed initial state, the corrected controls generate
\begin{subequations}\label{eq:rollout}\begin{align}
 \bar\bv_{i,k+1}&=\bar\bv_{i,k}+h\bar\ba_{i,k},\\
 \bar\bp_{i,k+1}&=\bar\bp_{i,k}+h\bar\bv_{i,k}+\tfrac12h^2\bar\ba_{i,k}.
\end{align}\end{subequations}
For $0\leq\tau\leq h$, the returned interval motion is
\begin{equation}
 \bar\bp_i(t_k+\tau)=\bar\bp_{i,k}+\tau\bar\bv_{i,k}
 +\tfrac12\tau^2\bar\ba_{i,k},\qquad
 \bar\bv_i(t_k+\tau)=\bar\bv_{i,k}+\tau\bar\ba_{i,k}.
 \label{eq:arc}
\end{equation}
It follows immediately from~\eqref{eq:rollout} that position and velocity agree at neighboring interval boundaries. Equation~\eqref{eq:correction} enforces both terminal moments, giving the prescribed final position and zero final velocity. The output therefore satisfies~\eqref{eq:continuous} on every interval, with continuous velocity and bounded, possibly discontinuous acceleration when the subsequent limit checks pass.

To relate reconstruction to numerical recovery, define
\begin{equation}
 D_i=\sum_{k=0}^{K-1}\norm{\mathbf E_{i,k}}_2^2,\qquad
 \mathbf e_i=\sum_{k=0}^{K-1}\mathbf E_{i,k},\qquad
 \norm{\mathbf e_i}_2\leq\sqrt{K D_i}.
 \label{eq:terminal-bound}
\end{equation}
The inequality follows from Cauchy--Schwarz. With exact candidate endpoint bounds and $\ba_{i,k}=(\bv_{i,k+1}-\bv_{i,k})/h$, the uncorrected terminal velocity is exact, and the remaining displacement mismatch is $\mathbf e_i$. For any coordinate $d$,~\eqref{eq:correction} therefore gives
\begin{equation}
 \norm{\bar{\mathbf a}_{i,d}-\mathbf a_{i,d}}_2
 \leq\left\|H^\top(HH^\top)^{-1}
 \begin{bmatrix}0\\1\end{bmatrix}\right\|_2\sqrt{K D_i}.
 \label{eq:correction-bound}
\end{equation}
For a fixed grid, the bound vanishes with $D_i$; the actual correction also depends on cancellation among interval defects. The integrated motion is checked against every obstacle and UAV pair, together with the speed, acceleration, and workspace bounds. These checks account for its displacement from the QP candidate while retaining the prescribed horizon and endpoint states.

\begin{algorithm2e}[htbp]
\caption{LIOM with hard-QP entry and acceptance-driven recovery}
\label{alg:liom}
\KwIn{Assigned starts and goals, map, $T_f$, motion limits, and budgets}
\KwOut{Accepted joint motion or planning failure}
Construct reference geometry and admissible scalar timing\;
\If{SFM generation fails}{Try static-guide geometry with the same timing rule\;}
\If{no admissible reference is available}{\Return{planning failure}\;}
Initialize $\rho_q,\rho_s,\varepsilon_q$ and the reference\;
\For{$\nu=0,\ldots,\nu_{\max}-1$}{
  \If{budget exhausted}{\textbf{break}\;}
  Construct spatial boxes, screen distant pairs, and construct relative boxes\;
  \If{boxes are invalid or exclude a fixed endpoint}{\textbf{break}\;}
  \If{$\nu=0$}{
    Try the bounded hard-QP entry~\eqref{eq:hardentry}\;
    \If{its reconstructed motion passes complete-interval checks}{\Return{accepted motion}\;}
  }
  Solve the relaxed QP~\eqref{eq:subproblem}\;
  \If{a finite candidate is available}{
    Evaluate $\bar D$, $\bar S$, and the local linear-constraint violation\;
    \If{numerical recovery and the candidate-readiness test hold}{
      Correct endpoint moments and integrate controls\;
      \If{complete-interval checks pass}{\Return{accepted motion}\;}
      Tighten $\varepsilon_q$ and increase $\rho_q$ within their limits\;
      Try exact kinematics and zero slack in the current boxes\;
      \If{the resulting reconstructed motion passes all checks}{\Return{accepted motion}\;}
    }
  }
  Update the two penalties independently using~\eqref{eq:weights}\;
  Update the reference from a usable candidate or the spatial-box centers\;
}
\Return{planning failure}\;
\end{algorithm2e}

\subsection{Computational Structure}
Acceleration elimination leaves $6N(K+1)$ state components before fixed endpoints are removed. With $M=|\mathcal A|$ active pair--node triples, the relaxed model adds $M$ slacks and $6M$ relative inequalities. Fixed-facet motion restrictions contribute $O(NK)$ rows. Sparse objective and constraint assembly therefore requires $O(NK+M)$ storage and work after the active set is known. Pair screening costs $O(N^2K)$ in the worst case. For a fixed expansion limit and step, testing $B$ obstacles during spatial-box construction costs $O(NKB)$, up to the fixed bisection count.

The complete-interval checker considers all pairs and obstacles, with a fixed-degree polynomial calculation per interval, giving $O((N^2+NB)K)$ work. QP solution time also depends on sparse factorization, numerical conditioning, and the required iterations. Geometric screening reduces actual coupling, but the centralized worst-case dependence remains quadratic in team size. Section~\ref{sec:scaling} measures both team and grid scaling.
